% Shared preamble of the HSEMHS mode documents (docs/physics/*.tex, docs/technical/*.tex).
\documentclass[10pt]{extarticle}
\usepackage[a4paper,margin=1.7cm]{geometry}
\usepackage[T1]{fontenc}\usepackage{lmodern}\usepackage{tabularx}
\usepackage{amsmath,amssymb,bm,enumitem,titlesec,xcolor,booktabs,array,tikz,graphicx,longtable}
\usepackage[hidelinks]{hyperref}
\titleformat{\section}{\large\bfseries\color{blue!50!black}}{\thesection.}{0.5em}{}
\titlespacing*{\section}{0pt}{9pt}{3pt}
\titleformat{\subsection}{\normalsize\bfseries}{\thesubsection}{0.5em}{}
\titlespacing*{\subsection}{0pt}{6pt}{2pt}
\setlist{nosep,leftmargin=1.3em}
\setlength{\parindent}{0pt}\setlength{\parskip}{4pt}
\renewcommand{\arraystretch}{1.15}
\newcommand{\B}{\mathbf B}\newcommand{\J}{\mathbf J}\newcommand{\F}{\mathbf F}\newcommand{\R}{\mathbf R}
\newcommand{\zh}{\hat{\mathbf z}}\newcommand{\Bo}{\mathbf B_{\rm obs}}\newcommand{\Lv}{\mathbf L}
\newcommand{\code}[1]{\texttt{#1}}
\newcolumntype{L}[1]{>{\raggedright\arraybackslash}p{#1}}
\newcolumntype{Y}{>{\raggedright\arraybackslash}X}
\newcommand{\doctitle}[2]{{\Large\bfseries #1}\\[2pt]{\small #2}\par\vspace{2pt}}
\pagestyle{plain}
\AtBeginDocument{\sloppy}

\begin{document}
\doctitle{column-hse-ltau: hydrostatic columns on an optical-depth grid}{Physics and method of \code{run.mode = "column-hse-ltau"} in HSEMHS. A self-contained description.}

\begin{abstract}
\noindent A Stokes inversion returns the temperature and the magnetic field of the solar atmosphere on surfaces of constant
optical depth, not at geometric heights. This mode turns such an atmosphere into a physical one: for every pixel column it
computes the gas pressure and the density on each optical-depth level and the geometric height at which the level lies,
assuming that each column is in hydrostatic equilibrium on its own. We derive the hydrostatic equation in optical depth,
show why it needs no externally supplied reference pressure and what replaces it, describe the closure at the top of the
column, the numerical integration, the safeguard against interpolation artefacts at sharp temperature jumps, and the rule
that places the columns relative to each other in height by equalising the total (gas plus magnetic) pressure. All
assumptions are listed.
\end{abstract}

\section{Introduction}
When we observe the Sun we receive light, and the light from a given depth is attenuated by the gas above it. A spectral
line inversion therefore naturally describes the atmosphere as a function of \textbf{optical depth} $\tau_{500}$ (the
optical depth in the continuum at 500\,nm): it gives $T$ and $\B$ at, for example, $\log_{10}\tau_{500}=-8.0,-7.9,\dots,+1.4$
in every pixel. Such an atmosphere has no pressure, no density and no geometry: we do not know how many kilograms of gas a
given $\tau$ interval contains, nor how many kilometres thick it is. Both follow from two physical laws, hydrostatic
equilibrium and the dependence of the opacity on the state of the gas.

\textbf{Input:} the $\log_{10}\tau_{500}$ levels, the horizontal coordinates $x$, $y$, and $T$, $B_x$, $B_y$, $B_z$ on every
level of every column; the gravity $g$. \textbf{Output:} on the same $\tau$ levels, $P$, $\rho$ and the height $z$, plus the
height of $\tau_{500}=1$ in every column. $T$ and $\B$ are copied unchanged.

\section{Notation}
\begin{tabularx}{\linewidth}{>{$}l<{$} Y l l}
\toprule
\text{Symbol} & Meaning & Unit & Role \\ \midrule
x,\ y & horizontal coordinates of the pixels, $x_i$, $y_j$, uniform spacing $\Delta x$, $\Delta y$ & m & given \\
\tau & optical depth at 500\,nm, counted from the top downward & -- & coordinate \\
\ell_k & $\log_{10}\tau_k$ of level $k=1\dots N_\tau$, increasing (top first), equally spaced ($\Delta\ell=0.1$) & -- & given \\
s & $\ln\tau=\ell\ln10$, the integration variable & -- & coordinate \\
T,\ X & temperature on the levels; $X=\ln T$ & K & given, fixed \\
\B & magnetic field on the levels, $(B_x,B_y,B_z)$ & T & given, fixed \\
P,\ u & gas pressure on the levels; $u=\ln(P/1\,{\rm Pa})$ & Pa, -- & unknown \\
z & geometric height of a level in its column (upward) & m & unknown \\
\rho & mass density $\rho(T,P)$ & kg\,m$^{-3}$ & derived \\
\chi & extinction coefficient at 500\,nm, $\chi_{500}(T,P)$ & m$^{-1}$ & derived \\
\kappa & extinction per unit mass, $\kappa=\chi/\rho$ & m$^2$\,kg$^{-1}$ & derived \\
q & $\rho/P$ & s$^2$\,m$^{-2}$ & derived \\
m & column mass above a point per unit area & kg\,m$^{-2}$ & derived \\
H & pressure scale height $P/(\rho g)$ & m & derived \\
P_{\rm tot} & total pressure $P+B^2/(2\mu_0)$ & Pa & derived \\
P_{\rm tot}^{\rm ref} & reference total pressure used to align the columns & Pa & derived \\
c_{\rm max} & largest allowed rise of $\ln P$ per decade of $\tau$ (\code{cap\_per\_dex}, default 50) & -- & parameter \\
g & gravitational acceleration, 274 & m\,s$^{-2}$ & parameter \\
\mu_0 & vacuum permeability $4\pi\times10^{-7}$ & T\,m\,A$^{-1}$ & constant \\
\bottomrule
\end{tabularx}

\section{Optical depth measures mass}\label{sec:taumass}
Going down by a height $dz$ (with $z$ upward, $dz<0$) the light from below has to cross an additional layer, and the optical
depth grows by
\begin{equation}
d\tau=-\chi\,dz .\label{eq:dtau}
\end{equation}
Write $\chi=\kappa\rho$: the extinction is the amount of gas, $\rho$, times how strongly each kilogram of it removes light,
$\kappa$. The mass of the column above a point grows by $dm=-\rho\,dz$ over the same step. Therefore
\begin{equation}
d\tau=\kappa\,dm .\label{eq:taumass}
\end{equation}
\emph{Optical depth is column mass weighted by the opacity.} For gas with a constant $\kappa$, $\tau=\kappa m$ exactly. This
is the key to the whole mode: an optical-depth scale carries information about how much mass lies above each level.

\section{Hydrostatic equilibrium in optical depth}\label{sec:hse}
\subsection{Derivation}
A thin horizontal slab of gas at rest carries its own weight by the pressure difference across it. With $z$ upward this is
the hydrostatic equation $dP/dz=-\rho g$ (the vertical component of $-\nabla P-\rho g\zh=0$). Since $dm=-\rho\,dz$ it can
also be written $dP=g\,dm$: the pressure at a point is the weight of the gas above it. Dividing by Eq.~\eqref{eq:taumass}:
\begin{equation}
\boxed{\;\frac{dP}{d\tau}=\frac{g}{\kappa(T,P)}\;}\label{eq:hsetau}
\end{equation}
Adding the optical depth $d\tau$ means adding the mass $d\tau/\kappa$, whose weight is $g\,d\tau/\kappa$. For constant
$\kappa$, $P=g\tau/\kappa$.

\subsection{The form that is integrated}
The levels are equally spaced in $\log_{10}\tau$, and $\tau$ spans ten decades. We therefore integrate in $s=\ln\tau$
($d\tau=\tau\,ds$) and in $u=\ln P$ ($dP=P\,du$). With $\kappa=\chi/\rho$:
\begin{equation}
\boxed{\;\frac{du}{ds}=\frac{g\,\tau}{\kappa\,P}=\frac{g\,\tau\,\rho}{\chi\,P}=\frac{g\,\tau\,q(T,P)}{\chi(T,P)}\;}\label{eq:duds}
\end{equation}
and, from Eq.~\eqref{eq:dtau}, the geometric thickness of the layers:
\begin{equation}
\boxed{\;\frac{dz}{ds}=-\frac{\tau}{\chi(T,P)}\;}\label{eq:dzds}
\end{equation}
Where the gas is opaque (large $\chi$) a step in $\ln\tau$ is a thin layer; where it is transparent, the same step is a thick
one. Equations~\eqref{eq:duds} and \eqref{eq:dzds} are two coupled ordinary differential equations along each column. $\rho$
and $\chi$ come from the equation-of-state table described next.

% Shared section of the physics documents: equation of state and 500 nm extinction (\input by each mode document).
\section{The equation of state and the opacity}\label{sec:eos}
Every force balance in this document involves the density $\rho$, and every optical depth involves the extinction
coefficient $\chi_{500}$. Neither is an unknown of its own: at given temperature and pressure both are fixed by the
microphysics of the gas. This section explains how.

\subsection{The ideal gas}
In the solar photosphere and chromosphere the particles interact only weakly, so the gas obeys the ideal-gas law. If $n$
is the number of free particles (atoms, ions, molecules and electrons) per cubic metre,
\begin{equation}
P=n\,k_BT .\label{eq:ideal}
\end{equation}
If the average mass of a free particle is $\bar\mu\,m_{\rm u}$ (with $m_{\rm u}$ the atomic mass unit and $\bar\mu$ the
\textbf{mean molecular weight}), the mass density is $\rho=n\,\bar\mu\,m_{\rm u}$, and therefore
\begin{equation}
\rho=\frac{\bar\mu\,m_{\rm u}}{k_B}\,\frac{P}{T}.\label{eq:rho_ideal}
\end{equation}
At fixed $\bar\mu$ and $T$ the density is proportional to the pressure.

\subsection{\texorpdfstring{Why $\bar\mu$ depends on $T$ and $P$}{Why the mean molecular weight depends on T and P}: ionisation and molecules}
For neutral solar gas (about 91\,\% hydrogen and 9\,\% helium atoms by number, plus a small fraction of heavier elements)
$\bar\mu\approx1.3$. When a hydrogen atom is ionised it becomes two free particles, a proton and an electron, of almost
the same total mass; in fully ionised gas $\bar\mu\approx0.6$. In the coolest layers hydrogen partly forms H$_2$ molecules,
which raises $\bar\mu$. The degree of ionisation of an element in stage $i$ is given by the \textbf{Saha equation}
\begin{equation}
\frac{n_{i+1}\,n_e}{n_i}=\frac{2U_{i+1}(T)}{U_i(T)}\left(\frac{2\pi m_ek_BT}{h^2}\right)^{3/2}e^{-E_i/k_BT},\label{eq:saha}
\end{equation}
with $E_i$ the ionisation energy (13.6\,eV for hydrogen), $U_i$ the partition functions, $m_e$ the electron mass and $h$
Planck's constant. Two properties matter for us. (i) Because of the exponential, ionisation switches on over a narrow
range of temperature. (ii) The electron density $n_e$ appears on the left: at higher density there are more electrons to
recombine with, and a higher temperature is needed for the same ionisation. Hence $\bar\mu=\bar\mu(T,P)$, and
Eq.~\eqref{eq:rho_ideal} is no longer a simple proportionality. The Saha equations of all elements are coupled through the
common $n_e$ and through charge conservation ($n_e$ equals the total charge of the ions), and the molecular equilibria
(H$_2$, H$_2^+$, H$^-$) are solved together with them. This assumes \textbf{local thermodynamic equilibrium} (LTE): the
level populations and ionisation stages are those of a gas in equilibrium at the local $T$.

\subsection{\texorpdfstring{The 500\,nm extinction}{The 500 nm extinction}}
Light of wavelength 500\,nm travelling a distance $ds$ loses the fraction $\chi_{500}\,ds$ of its intensity, by absorption
and by scattering. In the solar photosphere $\chi_{500}$ is dominated by the negative hydrogen ion H$^-$ (a hydrogen atom
with a second, weakly bound electron, binding energy 0.75\,eV), whose abundance is proportional to the electron density.
The electrons come mostly from easily ionised metals (Na, Mg, Si, Ca, Fe, \dots), so $\chi_{500}$ grows steeply with
temperature and with density. The code adds the continuum contributions of H (bound--free and free--free), H$^-$,
H$_2^+$, He, He$^-$, the metals C, N, Mg, Si, Ca and Fe, Rayleigh scattering on H, He and H$_2$, and Thomson scattering
on free electrons, each counted once. Spectral lines are not included: $\chi_{500}$ is a continuum extinction.

\subsection{The table and its interpolation}
Solving the coupled Saha, molecular and charge equations at every grid point at every step would be far too slow. They
are therefore solved once, with Wittmann's equation-of-state code for solar composition (with a dedicated hot branch above
a switch temperature $T_{\rm sw}$, where the original solver fails in the hot, low-pressure corner), on a grid of
$X=\ln T$ and $Y=\ln P$, and stored in a table. For every node the table holds the three functions
\[
\ln\rho(X,Y),\qquad \ln n_e(X,Y),\qquad \ln\chi_{500}(X,Y),
\]
together with their derivatives $\partial_X$, $\partial_Y$ and $\partial_X\partial_Y$. Between the nodes each function is
evaluated by \textbf{bicubic Hermite interpolation}: a cubic polynomial in each direction that matches the stored values
and derivatives at the four corners of the cell. The result is continuous with continuous first derivatives, which is
important for the Newton-type solvers of the magnetohydrostatic modes, and it provides the derivative
$\partial\ln\rho/\partial\ln P|_T$ that they need. The table used by all modes is \code{cache/eos\_table\_fullz.h5}.
Its range of $T$ and $P$ (its \emph{envelope}) covers the input atmospheres; a point outside the envelope is reported as
an equation-of-state failure, never extrapolated.

Two derived quantities appear repeatedly:
\begin{equation}
q\equiv\frac{\rho}{P}\quad[{\rm s^2\,m^{-2}}],\qquad H\equiv\frac{P}{\rho\,g}=\frac{1}{g\,q}=\frac{k_BT}{\bar\mu\,m_{\rm u}\,g}\quad[{\rm m}].
\label{eq:H}
\end{equation}
$H$ is the \textbf{pressure scale height}. For $T=6000$\,K and $\bar\mu=1.3$, $H\approx140$\,km.


\section{The boundary conditions}\label{sec:bc}
\subsection{Why no reference pressure is needed}\label{sec:noPref}
On a geometric height grid the hydrostatic equation $du/dz=-g\bar\mu m_{\rm u}/(k_BT)$ does not contain the pressure itself
when $\bar\mu$ is constant: if $P(z)$ is a solution, so is $c\,P(z)$ for any constant $c$. A height scale says nothing about
how much gas there is, and a reference pressure must be prescribed somewhere. On an optical-depth grid the situation is
different. The coordinate $\tau$ already counts the opacity-weighted mass (Eq.~\ref{eq:taumass}), and Eq.~\eqref{eq:hsetau}
has a natural starting point: at $\tau=0$ there is no gas above, so $P=0$. Integrating,
\begin{equation}
P(\tau)=\int_0^{\tau}\frac{g}{\kappa}\,d\tau' .\label{eq:Pint}
\end{equation}
Once $T(\tau)$ and the opacity law $\kappa(T,P)$ are known, the pressure follows; no value has to be supplied from outside.
Equivalently, Eq.~\eqref{eq:hsetau} is not invariant under $P\to cP$, because $\kappa$ depends on $P$ (the H$^-$ opacity
grows with the electron density).

\subsection{What replaces it: the top closure}
The grid does not start at $\tau=0$ but at $\tau_1=10^{-8}$. Above the top level we assume that $\kappa$ keeps its value
$\kappa_1$ and that $P(\tau=0)=0$. Equation~\eqref{eq:Pint} then gives
\begin{equation}
P_1=\frac{g\,\tau_1}{\kappa(T_1,P_1)}=\frac{g\,\tau_1\,\rho(T_1,P_1)}{\chi(T_1,P_1)} .\label{eq:top}
\end{equation}
$P_1$ appears on both sides, because $\kappa$ depends on $P$. It is found by a damped fixed-point iteration in $u$: start
from $u=\ln(100)$, compute $u_{\rm new}=\ln(g\tau_1\rho/\chi)$ at the current $u$, set $u\leftarrow\tfrac12(u+u_{\rm new})$,
and repeat until the change is below $10^{-10}$ (at most 200 times). The damping avoids oscillation.

How much does this assumption matter? At depth $\tau$ the pressure is $P_1$ plus the weight of the gas between $\tau_1$ and
$\tau$. If $\kappa$ were constant, $P_1/P(\tau)=\tau_1/\tau$, i.e.\ $10^{-8}$ at $\tau=1$. Any reasonable error in $\kappa$
above the top therefore affects only the uppermost few levels. The result is free of an external pressure, but it depends
on three modelling choices instead: this top closure, the treatment of sharp temperature jumps (Sect.~\ref{sec:cap}) and,
for the heights, the alignment of the columns (Sect.~\ref{sec:align}). It is also only as good as the opacity: if the
$\tau$ scale of the input was computed with a different opacity, the pressures differ in proportion.

\subsection{The height: a constant of integration per column}
Equation~\eqref{eq:dzds} fixes the \emph{distances} between the levels of a column, not where the column sits. Each column
is first integrated with $z=0$ at the top level and then shifted so that $z=0$ at $\log_{10}\tau=0$. Its absolute height is
fixed afterwards by the alignment rule of Sect.~\ref{sec:align}. So the reference that this mode needs is a reference
\emph{height}, not a reference pressure.

\section{Assumptions}\label{sec:assume}
\begin{enumerate}[label=\textbf{A\arabic*.},leftmargin=2.2em]
\item \textbf{Static gas, vertical balance only.} Each column is hydrostatic on its own; velocities and the horizontal force
balance between columns are ignored.
\item \textbf{No Lorentz force in the column balance.} The field enters only the alignment, through its magnetic pressure;
its tension is ignored.
\item \textbf{Consistent opacity.} The input $\tau_{500}$ scale is assumed to be the continuum optical depth at 500\,nm along
the vertical computed with the same LTE equation of state and opacity as the table.
\item \textbf{Top closure:} $\kappa$ constant above the top level and $P(\tau=0)=0$ (Eq.~\ref{eq:top}).
\item \textbf{Temperature between levels:} $\ln T$ varies linearly with $\ln\tau$ inside each interval.
\item \textbf{Limited pressure rise} of at most $c_{\rm max}\Delta\ell$ in $\ln P$ per interval (Sect.~\ref{sec:cap}).
\item \textbf{Alignment:} at one common height all columns have the same total pressure $P+B^2/2\mu_0$; below the deepest
level the total pressure is extrapolated with the gas-pressure scale height (Sect.~\ref{sec:align}).
\item \textbf{Gauge:} the mean height of $\tau_{500}=1$ is zero.
\item \textbf{Constant, vertical gravity} and plane-parallel geometry.
\item \textbf{Field attached to the $\tau$ levels:} the given field values stay on their levels wherever those end up in
space; nothing makes them divergence-free there.
\end{enumerate}

\section{Numerical integration of a column}\label{sec:num}
\subsection{Temperature inside an interval}
Between levels $k$ and $k+1$ ($s_k=\ell_k\ln10$, $s_{k+1}=\ell_{k+1}\ln10$),
\begin{equation}
X(s)=X_k+(X_{k+1}-X_k)\,\frac{s-s_k}{s_{k+1}-s_k} .
\end{equation}

\subsection{Heun's method with adaptive substeps}
Write Eqs.~\eqref{eq:duds}--\eqref{eq:dzds} as $du/ds=f(s,u)$, $dz/ds=h(s,u)$. The integration runs downward, from $s_k$ to
$s_{k+1}$, in substeps of length $\delta$. One substep from $(s,u,z)$ is Heun's predictor--corrector method (the
trapezoidal rule with an Euler predictor):
\begin{equation}
\begin{aligned}
&f_1=f(s,u),\quad h_1=h(s,u),\qquad \delta=\min\!\left(s_{k+1}-s,\ \frac{0.05}{|f_1|}\right),\\
&\tilde u=u+\delta f_1\ \ \text{(predictor)},\qquad f_2=f(s+\delta,\tilde u),\quad h_2=h(s+\delta,\tilde u),\\
&u\leftarrow u+\tfrac{\delta}{2}(f_1+f_2),\qquad z\leftarrow z+\tfrac{\delta}{2}(h_1+h_2),\qquad s\leftarrow s+\delta .
\end{aligned}\label{eq:heun}
\end{equation}
The substep length is chosen so that $\ln P$ changes by at most about 0.05 per substep (a 5\,\% change in pressure), so
intervals where the pressure rises steeply receive many substeps. The method is second-order accurate: the error falls with
the square of the substep length.

\subsection{Hot gas over cool gas: the limit on the pressure rise}\label{sec:cap}
The 500\,nm opacity rises very steeply with temperature (Sect.~\ref{sec:eos}), so cool gas is far more transparent than hot
gas. Suppose that within one grid interval a hot layer (8000\,K) lies directly on a cool one (4000\,K). With the linear
interpolation of $\ln T$, a whole tenth of a decade in $\tau$ is then filled with nearly transparent gas. By
Eq.~\eqref{eq:hsetau}, collecting that much optical depth from transparent gas requires an enormous mass, and the pressure
rises by many orders of magnitude within the interval. We regard this as mainly a consequence of representing a sharp
temperature change by smooth interpolation on a coarse grid, although the true structure between two levels is unknown.
Two limits therefore apply within each interval:
\begin{itemize}
\item $\ln P$ may rise by at most $c_{\rm max}\,\Delta\ell=50\times0.1=5$ across the interval (a factor of about 150 in
pressure), well above the rises of ordinary stratifications;
\item $\ln P$ may not exceed the upper edge of the table minus 1, and the slope must remain finite.
\end{itemize}
If a substep would violate either, the pressure at the lower level is set to the limited value
$u_{k+1}=\min(u_k+c_{\rm max}\Delta\ell,\ u_{\max})$, and the integration of that interval stops there. In such an interval
the column is \emph{not} hydrostatic, and the part of its geometric thickness that was not yet integrated is not added to the
heights below (a known shortcoming). Below it, the hydrostatic integration continues from the limited pressure. The number
of limited intervals is reported; results in those columns should be treated with care.

\subsection{The result of a column}
After the last interval each column has $u_k$ and $z_k$ on all levels, with $z_1=0$. The heights are shifted so that $z=0$ at
$\ell=0$, the value at $\ell=0$ being interpolated linearly in $\ell$ (or extrapolated from the end interval). If any level
of any column ends with a non-finite pressure (outside the table), the run stops with an error.

\section{Where the columns sit: total-pressure alignment}\label{sec:align}
\subsection{Magnetic pressure}
After Sect.~\ref{sec:num} every column has a correct internal structure but an arbitrary vertical position. To place the
columns relative to each other we use the horizontal force balance in its simplest form. The magnetic (Lorentz) force can be
rewritten with a vector identity:
\begin{equation}
\J\times\B=\frac{(\nabla\times\B)\times\B}{\mu_0}
=-\nabla\!\left(\frac{B^2}{2\mu_0}\right)+\frac{(\B\cdot\nabla)\B}{\mu_0}.
\end{equation}
The first term acts like an additional pressure, the \textbf{magnetic pressure} $B^2/2\mu_0$; the second is a tension along
curved field lines.

\subsection{The alignment rule}
If the field is nearly vertical and weakly curved, its tension is small, and the horizontal force balance between
neighbouring columns at the same height reduces to equal total pressure:
\begin{equation}
\frac{\partial}{\partial x}\left(P+\frac{B^2}{2\mu_0}\right)=\frac{\partial}{\partial y}\left(P+\frac{B^2}{2\mu_0}\right)=0
\qquad\text{(tension neglected)} .
\end{equation}
A column with a higher total pressure would push its neighbour aside. The mode imposes this at \emph{one} height only:
\begin{enumerate}
\item In every column compute $\ln P_{\rm tot,k}=\ln\!\big(P_k+B_k^2/2\mu_0\big)$ on all levels.
\item The reference is the median over all columns of the total pressure on the deepest level,
$P_{\rm tot}^{\rm ref}=\operatorname{median}_{i,j}P_{{\rm tot},N_\tau}$.
\item In each column find the height $z_r$ at which $P_{\rm tot}=P_{\rm tot}^{\rm ref}$: go up from the deepest level to the
first interval in which $\ln P_{\rm tot}$ drops below $\ln P_{\rm tot}^{\rm ref}$, and interpolate linearly in $z$ between its
two levels.
\item For a column whose deepest total pressure is already below the reference, the crossing lies below the grid. It is
extrapolated with the gas-pressure scale height of the deepest level,
$z_r=z_{N_\tau}-\big(\ln P_{\rm tot}^{\rm ref}-\ln P_{{\rm tot},N_\tau}\big)\,H_{N_\tau}$ with $H=P/(\rho g)$. This is an
approximation (the total pressure does not fall with exactly the gas scale height); the number of such columns is reported.
\item Shift the column: $z_k\leftarrow z_k-z_r$. Now $P_{\rm tot}=P_{\rm tot}^{\rm ref}$ at $z=0$ in every column.
\end{enumerate}
With \code{align = "none"} this step is skipped and every column keeps $z=0$ at $\tau_{500}=1$.

\subsection{Consequence: the Wilson depression}
Under this rule, a column with a strong field carries part of the total pressure as magnetic pressure, so it needs less gas
pressure at the reference height. Less gas pressure means lower density and lower opacity, so light escapes from deeper:
its $\tau$ surfaces lie lower. This is the \textbf{Wilson depression} seen in sunspots observed near the limb. It is an
expected tendency, not a guarantee: the outcome for a given column also depends on its temperature structure and on the
approximations above.

\subsection{The gauge}
Finally all columns are shifted together by the same amount so that the mean over all columns of the height of
$\tau_{500}=1$ is zero. This is only a choice of origin.

\section{Iteration and acceptance criteria}
The mode contains no global iteration. Its convergence criteria are local:
\begin{itemize}
\item the top pressure: change of $u$ below $10^{-10}$ in the fixed-point iteration of Eq.~\eqref{eq:top};
\item the substep length: change of $\ln P$ at most 0.05 per Heun substep;
\item the pressure rise limit and the table edge (Sect.~\ref{sec:cap}), counted and reported.
\end{itemize}
Every column is independent until the alignment, which needs one global number (the median reference total pressure)
and the gauge, which needs one global mean.

\section{Three dimensions: what the result does and does not satisfy}
\begin{itemize}
\item \textbf{Vertical:} every column obeys $dP/d\tau=g/\kappa$ (hydrostatic equilibrium), except in limited intervals.
\item \textbf{Horizontal:} the total pressure of all columns is equal at $z=0$ only. At other heights neighbouring columns
generally have different total pressures, so a horizontal force remains; the field's tension is not included at all. This
mode is therefore not a three-dimensional equilibrium. Imposing horizontal force balance at all heights, together with the
full Lorentz force, requires moving the $\tau$ surfaces in height and adjusting the pressure in a coupled way, which is what
the mode \code{mhs-poisson-ltau} does starting from this result.
\item \textbf{The field:} the given $\B$ values stay on their $\tau$ levels. Because the levels move to new heights, the field
as it ends up in space need not satisfy $\nabla\cdot\B=0$; this mode does not use derivatives of $\B$ and is not affected,
but any later use of the field should keep it in mind.
\end{itemize}

\section{Summary of the algorithm}
\begin{enumerate}
\item Read $\ell_k$, $x$, $y$, $T$, $\B$; load the equation-of-state table.
\item For every column independently: find $P_1=g\tau_1\rho/\chi$ by damped fixed-point iteration; integrate $du/ds=g\tau
q/\chi$ and $dz/ds=-\tau/\chi$ downward with Heun substeps ($\Delta\ln P\le0.05$), $\ln T$ linear in $\ln\tau$, limiting the
rise of $\ln P$ to $c_{\rm max}\Delta\ell$ per interval; shift so that $z=0$ at $\tau=1$.
\item Compute the median total pressure on the deepest level; shift every column so that its total pressure equals it at
$z=0$ (extrapolating below the grid where needed).
\item Shift all columns so that the mean height of $\tau_{500}=1$ is zero.
\item Evaluate $\rho$ from the table; write $z$, $P$, $\rho$, $T$, $\B$ and $z(\tau_{500}=1)$.
\end{enumerate}
\end{document}
